\documentclass[10pt]{article}
\usepackage[utf8]{inputenc}
\usepackage[T1]{fontenc}
\usepackage{amsmath,amssymb,adjustbox}
\usepackage[margin=2cm]{geometry}
\begin{document}
\begin{center}\large Reducción y resolución formal\end{center}
\textbf{Modelo:} R\_plus\_alpha\_times\_RX\par
\[\adjustbox{max width=\linewidth}{$\displaystyle \text{Lagrangiano:}\qquad L=g{}^{a\,c}\,g{}^{b\,d}\,R{}_{a\,b\,c\,d} + \alpha\,g{}^{a\,c}\,g{}^{b\,d}\,R{}_{a\,b\,c\,d}\,g{}^{a\,b}\,u{}_{a}\,u{}_{b}$}\]
\textbf{Ansatz:} draft4\_circular\_specialized\par
\[\adjustbox{max width=\linewidth}{$\displaystyle \text{Especialización de }\phi:\quad \phi=q\,\varphi$}\]
\textbf{Estado global:} hay familias verificadas y ramas pendientes\par
\[\adjustbox{max width=\linewidth}{$\displaystyle \text{Dominio de trabajo:}\qquad f\!\left(r\right)\neq0,\quad r\neq0$}\]
\textbf{Supuestos declarados:} r>0; f(r)!=0\par
\section*{Ecuaciones combinadas}
\[\adjustbox{max width=\linewidth}{$\displaystyle E^{\tau}_{\tau}-E^{r}_{r}=0$}\]
\[\adjustbox{max width=\linewidth}{$\displaystyle E^{\tau}_{\tau}-E^{\varphi}_{\varphi}=0$}\]
\[\adjustbox{max width=\linewidth}{$\displaystyle E^{r}_{r}-E^{\varphi}_{\varphi}=0$}\]
\section*{Combinaciones proyectadas y reducidas}
\[\adjustbox{max width=\linewidth}{$\displaystyle E^{\tau}_{\tau}-E^{r}_{r}=\frac{6\,\alpha\,{q}^{2}\,f\!\left(r\right)}{{r}^{4}}=0$}\]
\[\adjustbox{max width=\linewidth}{$\displaystyle E^{\tau}_{\tau}-E^{\varphi}_{\varphi}=\frac{1}{2\,{r}^{4}}\,\left(f'\!\left(r\right)\,{r}^{3} - f''\!\left(r\right)\,{r}^{4} + \alpha\,{q}^{2}\,\left(-4\,f\!\left(r\right) + f''\!\left(r\right)\,{r}^{2} + 7\,r\,f'\!\left(r\right)\right)\right)=0$}\]
\[\adjustbox{max width=\linewidth}{$\displaystyle E^{r}_{r}-E^{\varphi}_{\varphi}=\frac{1}{2\,{r}^{4}}\,\left(f'\!\left(r\right)\,{r}^{3} - f''\!\left(r\right)\,{r}^{4} + \alpha\,{q}^{2}\,\left(-16\,f\!\left(r\right) + f''\!\left(r\right)\,{r}^{2} + 7\,r\,f'\!\left(r\right)\right)\right)=0$}\]
\section*{Ecuaciones absolutas y escalares}
\[\adjustbox{max width=\linewidth}{$\displaystyle E^{\tau}_{\tau}=-\frac{\left(-f'\!\left(r\right)\,{r}^{3} + \alpha\,{q}^{2}\,\left(-8\,f\!\left(r\right) + r\,f'\!\left(r\right)\right)\right)}{2\,{r}^{4}}=0$}\]
\[\adjustbox{max width=\linewidth}{$\displaystyle E_\phi=0=0$}\]
Componentes fuera de la diagonal: $0=0$.\par
\section*{Familias verificadas}
Familia 1: verificada en todas las ecuaciones, en el dominio declarado.\par
\[\adjustbox{max width=\linewidth}{$\displaystyle \alpha=0$}\]
\[\adjustbox{max width=\linewidth}{$\displaystyle q=0$}\]
\[\adjustbox{max width=\linewidth}{$\displaystyle f\!\left(r\right)=C_{1}$}\]
\textbf{Método:} búsqueda de soluciones constantes\par
Parámetros libres en esta familia: $C_{1}$.\par
\[\adjustbox{max width=\linewidth}{$\displaystyle \text{Restricciones no nulas:}\qquad C_{1}\neq0,\quad r\neq0,\quad C_{1}\neq0$}\]
Supuestos: r>0; f(r)!=0\par
Familia 2: verificada en todas las ecuaciones, en el dominio declarado.\par
\[\adjustbox{max width=\linewidth}{$\displaystyle \alpha=0$}\]
\[\adjustbox{max width=\linewidth}{$\displaystyle f\!\left(r\right)=C_{1}$}\]
\textbf{Método:} búsqueda de soluciones constantes\par
Parámetros libres en esta familia: $C_{1},\; q$.\par
\[\adjustbox{max width=\linewidth}{$\displaystyle \text{Restricciones no nulas:}\qquad C_{1}\neq0,\quad r\neq0,\quad q\neq0,\quad C_{1}\neq0$}\]
Supuestos: r>0; f(r)!=0\par
Familia 3: verificada en todas las ecuaciones, en el dominio declarado.\par
\[\adjustbox{max width=\linewidth}{$\displaystyle q=0$}\]
\[\adjustbox{max width=\linewidth}{$\displaystyle f\!\left(r\right)=C_{1}$}\]
\textbf{Método:} búsqueda de soluciones constantes\par
Parámetros libres en esta familia: $C_{1},\; \alpha$.\par
\[\adjustbox{max width=\linewidth}{$\displaystyle \text{Restricciones no nulas:}\qquad C_{1}\neq0,\quad r\neq0,\quad \alpha\neq0,\quad C_{1}\neq0$}\]
Supuestos: r>0; f(r)!=0\par
\section*{Auditoría de ramas no aceptadas}
Candidatos conservados solo en results.json: rejected=32, undetermined=2, partial=8.\par
Sistema: DAE; orden máximo: 2.\par
Salida formal conservada en Wolfram: I\par
Completitud: no demostrada. Una búsqueda simbólica acotada no certifica que no existan otras ramas diferenciales o singulares.\par
\end{document}